\IfFileExists{stacks-project.cls}{%
\documentclass{stacks-project}
}{%
\documentclass{amsart}
}

% For dealing with references we use the comment environment
\usepackage{verbatim}
\newenvironment{reference}{\comment}{\endcomment}
%\newenvironment{reference}{}{}
\newenvironment{slogan}{\comment}{\endcomment}
\newenvironment{history}{\comment}{\endcomment}

% For commutative diagrams we use Xy-pic
\usepackage[all]{xy}

% We use 2cell for 2-commutative diagrams.
\xyoption{2cell}
\UseAllTwocells

% We use multicol for the list of chapters between chapters
\usepackage{multicol}

% This is generally recommended for better output
\usepackage{lmodern}
\usepackage[T1]{fontenc}

% For cross-file-references
\usepackage{xr-hyper}

% Package for hypertext links:
\usepackage{hyperref}

% For any local file, say "hello.tex" you want to link to please
% use \externaldocument[hello-]{hello}
\externaldocument[introduction-]{introduction}
\externaldocument[conventions-]{conventions}
\externaldocument[sets-]{sets}
\externaldocument[categories-]{categories}
\externaldocument[topology-]{topology}
\externaldocument[sheaves-]{sheaves}
\externaldocument[sites-]{sites}
\externaldocument[stacks-]{stacks}
\externaldocument[fields-]{fields}
\externaldocument[algebra-]{algebra}
\externaldocument[brauer-]{brauer}
\externaldocument[homology-]{homology}
\externaldocument[derived-]{derived}
\externaldocument[simplicial-]{simplicial}
\externaldocument[more-algebra-]{more-algebra}
\externaldocument[smoothing-]{smoothing}
\externaldocument[modules-]{modules}
\externaldocument[sites-modules-]{sites-modules}
\externaldocument[injectives-]{injectives}
\externaldocument[cohomology-]{cohomology}
\externaldocument[sites-cohomology-]{sites-cohomology}
\externaldocument[dga-]{dga}
\externaldocument[dpa-]{dpa}
\externaldocument[sdga-]{sdga}
\externaldocument[hypercovering-]{hypercovering}
\externaldocument[schemes-]{schemes}
\externaldocument[constructions-]{constructions}
\externaldocument[properties-]{properties}
\externaldocument[morphisms-]{morphisms}
\externaldocument[coherent-]{coherent}
\externaldocument[divisors-]{divisors}
\externaldocument[limits-]{limits}
\externaldocument[varieties-]{varieties}
\externaldocument[topologies-]{topologies}
\externaldocument[descent-]{descent}
\externaldocument[perfect-]{perfect}
\externaldocument[more-morphisms-]{more-morphisms}
\externaldocument[flat-]{flat}
\externaldocument[groupoids-]{groupoids}
\externaldocument[more-groupoids-]{more-groupoids}
\externaldocument[etale-]{etale}
\externaldocument[chow-]{chow}
\externaldocument[intersection-]{intersection}
\externaldocument[pic-]{pic}
\externaldocument[weil-]{weil}
\externaldocument[adequate-]{adequate}
\externaldocument[dualizing-]{dualizing}
\externaldocument[duality-]{duality}
\externaldocument[discriminant-]{discriminant}
\externaldocument[derham-]{derham}
\externaldocument[local-cohomology-]{local-cohomology}
\externaldocument[algebraization-]{algebraization}
\externaldocument[curves-]{curves}
\externaldocument[resolve-]{resolve}
\externaldocument[models-]{models}
\externaldocument[functors-]{functors}
\externaldocument[equiv-]{equiv}
\externaldocument[pione-]{pione}
\externaldocument[etale-cohomology-]{etale-cohomology}
\externaldocument[proetale-]{proetale}
\externaldocument[relative-cycles-]{relative-cycles}
\externaldocument[more-etale-]{more-etale}
\externaldocument[trace-]{trace}
\externaldocument[crystalline-]{crystalline}
\externaldocument[spaces-]{spaces}
\externaldocument[spaces-properties-]{spaces-properties}
\externaldocument[spaces-morphisms-]{spaces-morphisms}
\externaldocument[decent-spaces-]{decent-spaces}
\externaldocument[spaces-cohomology-]{spaces-cohomology}
\externaldocument[spaces-limits-]{spaces-limits}
\externaldocument[spaces-divisors-]{spaces-divisors}
\externaldocument[spaces-over-fields-]{spaces-over-fields}
\externaldocument[spaces-topologies-]{spaces-topologies}
\externaldocument[spaces-descent-]{spaces-descent}
\externaldocument[spaces-perfect-]{spaces-perfect}
\externaldocument[spaces-more-morphisms-]{spaces-more-morphisms}
\externaldocument[spaces-flat-]{spaces-flat}
\externaldocument[spaces-groupoids-]{spaces-groupoids}
\externaldocument[spaces-more-groupoids-]{spaces-more-groupoids}
\externaldocument[bootstrap-]{bootstrap}
\externaldocument[spaces-pushouts-]{spaces-pushouts}
\externaldocument[spaces-chow-]{spaces-chow}
\externaldocument[groupoids-quotients-]{groupoids-quotients}
\externaldocument[spaces-more-cohomology-]{spaces-more-cohomology}
\externaldocument[spaces-simplicial-]{spaces-simplicial}
\externaldocument[spaces-duality-]{spaces-duality}
\externaldocument[formal-spaces-]{formal-spaces}
\externaldocument[restricted-]{restricted}
\externaldocument[spaces-resolve-]{spaces-resolve}
\externaldocument[formal-defos-]{formal-defos}
\externaldocument[defos-]{defos}
\externaldocument[cotangent-]{cotangent}
\externaldocument[examples-defos-]{examples-defos}
\externaldocument[algebraic-]{algebraic}
\externaldocument[examples-stacks-]{examples-stacks}
\externaldocument[stacks-sheaves-]{stacks-sheaves}
\externaldocument[criteria-]{criteria}
\externaldocument[artin-]{artin}
\externaldocument[quot-]{quot}
\externaldocument[stacks-properties-]{stacks-properties}
\externaldocument[stacks-morphisms-]{stacks-morphisms}
\externaldocument[stacks-limits-]{stacks-limits}
\externaldocument[stacks-cohomology-]{stacks-cohomology}
\externaldocument[stacks-perfect-]{stacks-perfect}
\externaldocument[stacks-introduction-]{stacks-introduction}
\externaldocument[stacks-more-morphisms-]{stacks-more-morphisms}
\externaldocument[stacks-geometry-]{stacks-geometry}
\externaldocument[moduli-]{moduli}
\externaldocument[moduli-curves-]{moduli-curves}
\externaldocument[examples-]{examples}
\externaldocument[exercises-]{exercises}
\externaldocument[guide-]{guide}
\externaldocument[desirables-]{desirables}
\externaldocument[coding-]{coding}
\externaldocument[obsolete-]{obsolete}
\externaldocument[fdl-]{fdl}
\externaldocument[index-]{index}

% Theorem environments.
%
\theoremstyle{plain}
\newtheorem{theorem}[subsection]{Theorem}
\newtheorem{proposition}[subsection]{Proposition}
\newtheorem{lemma}[subsection]{Lemma}

\theoremstyle{definition}
\newtheorem{definition}[subsection]{Definition}
\newtheorem{example}[subsection]{Example}
\newtheorem{exercise}[subsection]{Exercise}
\newtheorem{situation}[subsection]{Situation}

\theoremstyle{remark}
\newtheorem{remark}[subsection]{Remark}
\newtheorem{remarks}[subsection]{Remarks}

\numberwithin{equation}{subsection}

% Macros
%
\def\lim{\mathop{\mathrm{lim}}\nolimits}
\def\colim{\mathop{\mathrm{colim}}\nolimits}
\def\Spec{\mathop{\mathrm{Spec}}}
\def\Hom{\mathop{\mathrm{Hom}}\nolimits}
\def\Ext{\mathop{\mathrm{Ext}}\nolimits}
\def\SheafHom{\mathop{\mathcal{H}\!\mathit{om}}\nolimits}
\def\SheafExt{\mathop{\mathcal{E}\!\mathit{xt}}\nolimits}
\def\Sch{\mathit{Sch}}
\def\Mor{\mathop{\mathrm{Mor}}\nolimits}
\def\Ob{\mathop{\mathrm{Ob}}\nolimits}
\def\Sh{\mathop{\mathit{Sh}}\nolimits}
\def\NL{\mathop{N\!L}\nolimits}
\def\CH{\mathop{\mathrm{CH}}\nolimits}
\def\proetale{{pro\text{-}\acute{e}tale}}
\def\etale{{\acute{e}tale}}
\def\QCoh{\mathit{QCoh}}
\def\Ker{\mathop{\mathrm{Ker}}}
\def\Im{\mathop{\mathrm{Im}}}
\def\Coker{\mathop{\mathrm{Coker}}}
\def\Coim{\mathop{\mathrm{Coim}}}

% Boxtimes
%
\DeclareMathSymbol{\boxtimes}{\mathbin}{AMSa}{"02}

%
% Macros for moduli stacks/spaces
%
\def\QCohstack{\mathcal{QC}\!\mathit{oh}}
\def\Cohstack{\mathcal{C}\!\mathit{oh}}
\def\Spacesstack{\mathcal{S}\!\mathit{paces}}
\def\Quotfunctor{\mathrm{Quot}}
\def\Hilbfunctor{\mathrm{Hilb}}
\def\Curvesstack{\mathcal{C}\!\mathit{urves}}
\def\Polarizedstack{\mathcal{P}\!\mathit{olarized}}
\def\Complexesstack{\mathcal{C}\!\mathit{omplexes}}
% \Pic is the operator that assigns to X its picard group, usage \Pic(X)
% \Picardstack_{X/B} denotes the Picard stack of X over B
% \Picardfunctor_{X/B} denotes the Picard functor of X over B
\def\Pic{\mathop{\mathrm{Pic}}\nolimits}
\def\Picardstack{\mathcal{P}\!\mathit{ic}}
\def\Picardfunctor{\mathrm{Pic}}
\def\Deformationcategory{\mathcal{D}\!\mathit{ef}}
\usepackage{cleveref}
\crefname{originalreference}{original reference}{original references}
\Crefname{originalreference}{Original reference}{Original references}
\makeatletter
\newlabel{categories-lemma-composition-cartesian}{{02XL}{source}{}{sourcepointer-1}{}}
\newlabel{categories-lemma-composition-cartesian@cref}{{[originalreference][0][]02XL}{[1][0][]source}{}{}{}}
\makeatother
\crefname{originalreference}{original reference}{original references}
\Crefname{originalreference}{Original reference}{Original references}
\makeatletter
\newlabel{categories-lemma-equality-morphisms-right-localization}{{04VJ}{source}{}{sourcepointer-2}{}}
\newlabel{categories-lemma-equality-morphisms-right-localization@cref}{{[originalreference][0][]04VJ}{[1][0][]source}{}{}{}}
\makeatother
\crefname{originalreference}{original reference}{original references}
\Crefname{originalreference}{Original reference}{Original references}
\makeatletter
\newlabel{categories-lemma-morphisms-right-localization}{{04VI}{source}{}{sourcepointer-3}{}}
\newlabel{categories-lemma-morphisms-right-localization@cref}{{[originalreference][0][]04VI}{[1][0][]source}{}{}{}}
\makeatother
\crefname{originalreference}{original reference}{original references}
\Crefname{originalreference}{Original reference}{Original references}
\makeatletter
\newlabel{categories-lemma-properties-right-localization}{{04VK}{source}{}{sourcepointer-4}{}}
\newlabel{categories-lemma-properties-right-localization@cref}{{[originalreference][0][]04VK}{[1][0][]source}{}{}{}}
\makeatother
\crefname{originalreference}{original reference}{original references}
\Crefname{originalreference}{Original reference}{Original references}
\makeatletter
\newlabel{categories-lemma-right-localization}{{04VH}{source}{}{sourcepointer-5}{}}
\newlabel{categories-lemma-right-localization@cref}{{[originalreference][0][]04VH}{[1][0][]source}{}{}{}}
\makeatother
\crefname{originalreference}{original reference}{original references}
\Crefname{originalreference}{Original reference}{Original references}
\makeatletter
\newlabel{categories-section-localization}{{04VB}{source}{}{sourcepointer-6}{}}
\newlabel{categories-section-localization@cref}{{[originalreference][0][]04VB}{[1][0][]source}{}{}{}}
\makeatother
\crefname{originallemma}{original lemma}{original lemmas}
\Crefname{originallemma}{Original lemma}{Original lemmas}
\makeatletter
\newlabel{lemma-right-multiplicative-system}{{04WF}{source}{}{sourcepointer-7}{}}
\newlabel{lemma-right-multiplicative-system@cref}{{[originallemma][0][]04WF}{[1][0][]source}{}{}{}}
\makeatother
\title{The right-localized pullback of a fibred category is fibred}
\author{The Stacks Project Authors}
\date{stacks-project:stacks.tex:local-snapshot-20261012; GNU FDL 1.2 or later}
\begin{document}
\maketitle
\noindent\textit{Editorial note.} The proof explicitly omits the converse sub-claim (that a strongly cartesian fraction has strongly cartesian fibre component), stating it is not needed for this lemma's conclusion; the fibredness conclusion is fully carried by the displayed arguments.\par\medskip
\noindent\textit{Allowed standard prerequisites (editorial).}\begin{itemize}
\item Right localization of categories. categories-section-localization, categories-lemma-right-localization, categories-lemma-properties-right-localization (cross-chapter, named). Describe u\_pS as the localization of u\_\{pp\}S at R and give the canonical functor p\_p extending p\_\{pp\}.
\item Morphisms and equality in the localization. categories-lemma-morphisms-right-localization, categories-lemma-equality-morphisms-right-localization (cross-chapter, named). Represent morphisms of u\_pS as fractions and decide equality after further localization; used in the uniqueness part.
\item Composition of strongly cartesian morphisms. categories-lemma-composition-cartesian (cross-chapter, named). Reduces strong cartesianness of a fraction to that of its numerator.
\end{itemize}\medskip
\ifdefined\hypertarget\hypertarget{sourcepointer-1}{}\fi
\noindent\textit{Source reference.} Original reference 02XL: focused-seeds/geometry/S9-stk-fibred-category-pullback/ref-tags.txt, source lines 5--5. Navigation pointer only: the counted proof cites this statement by \textbackslash{}ref as a standing input; the deduction the counted proof has to make is carried out inside the counted span, and no part of that deduction is delegated to this pointer.\par\medskip
\ifdefined\hypertarget\hypertarget{sourcepointer-2}{}\fi
\noindent\textit{Source reference.} Original reference 04VJ: focused-seeds/geometry/S9-stk-fibred-category-pullback/ref-tags.txt, source lines 6--6. Navigation pointer only: the counted proof cites this statement by \textbackslash{}ref as a standing input; the deduction the counted proof has to make is carried out inside the counted span, and no part of that deduction is delegated to this pointer.\par\medskip
\ifdefined\hypertarget\hypertarget{sourcepointer-3}{}\fi
\noindent\textit{Source reference.} Original reference 04VI: focused-seeds/geometry/S9-stk-fibred-category-pullback/ref-tags.txt, source lines 7--7. Navigation pointer only: the counted proof cites this statement by \textbackslash{}ref as a standing input; the deduction the counted proof has to make is carried out inside the counted span, and no part of that deduction is delegated to this pointer.\par\medskip
\ifdefined\hypertarget\hypertarget{sourcepointer-4}{}\fi
\noindent\textit{Source reference.} Original reference 04VK: focused-seeds/geometry/S9-stk-fibred-category-pullback/ref-tags.txt, source lines 8--8. Navigation pointer only: the counted proof cites this statement by \textbackslash{}ref as a standing input; the deduction the counted proof has to make is carried out inside the counted span, and no part of that deduction is delegated to this pointer.\par\medskip
\ifdefined\hypertarget\hypertarget{sourcepointer-5}{}\fi
\noindent\textit{Source reference.} Original reference 04VH: focused-seeds/geometry/S9-stk-fibred-category-pullback/ref-tags.txt, source lines 9--9. Navigation pointer only: the counted proof cites this statement by \textbackslash{}ref as a standing input; the deduction the counted proof has to make is carried out inside the counted span, and no part of that deduction is delegated to this pointer.\par\medskip
\ifdefined\hypertarget\hypertarget{sourcepointer-6}{}\fi
\noindent\textit{Source reference.} Original reference 04VB: focused-seeds/geometry/S9-stk-fibred-category-pullback/ref-tags.txt, source lines 10--10. Navigation pointer only: the counted proof cites this statement by \textbackslash{}ref as a standing input; the deduction the counted proof has to make is carried out inside the counted span, and no part of that deduction is delegated to this pointer.\par\medskip
\ifdefined\hypertarget\hypertarget{sourcepointer-7}{}\fi
\noindent\textit{Source reference.} Original lemma 04WF: sources/stacks-stacks/pinned-original-source.tex, source lines 3078--3078. Navigation pointer only: the counted proof cites this statement by \textbackslash{}ref as a standing input; the deduction the counted proof has to make is carried out inside the counted span, and no part of that deduction is delegated to this pointer.\par\medskip
\begin{lemma}
\label{lemma-fibred-category-pullback}
With notation and assumptions as in
Lemma \ref{lemma-right-multiplicative-system}.
Set $u_p\mathcal{S} = R^{-1}u_{pp}\mathcal{S}$, see
Categories, Section \ref{categories-section-localization}.
Then $u_p\mathcal{S}$ is a fibred category over $\mathcal{D}$.
\end{lemma}

\begin{proof}
We use the description of $u_p\mathcal{S}$ given just above
Categories, Lemma \ref{categories-lemma-right-localization}.
Note that the functor $p_{pp} : u_{pp}\mathcal{S} \to \mathcal{D}$
transforms every element of $R$ to an identity morphism.
Hence by
Categories, Lemma \ref{categories-lemma-properties-right-localization}
we obtain a canonical functor $p_p : u_p\mathcal{S} \to \mathcal{D}$
extending the given functor. This is how we think of
$u_p\mathcal{S}$ as a category over $\mathcal{D}$.

\medskip\noindent
First we want to characterize the $\mathcal{D}$-strongly cartesian morphisms
in $u_p\mathcal{S}$.
A morphism $f : X \to Y$ of $u_p\mathcal{S}$ is the equivalence class of
a pair $(f' : X' \to Y, r : X' \to X)$ with $r \in R$.
In fact, in $u_p\mathcal{S}$ we have $f = (f', 1) \circ (r, 1)^{-1}$
with obvious notation.
Note that an isomorphism is always strongly cartesian, as are
compositions of strongly cartesian morphisms, see
Categories, Lemma \ref{categories-lemma-composition-cartesian}.
Hence $f$ is strongly cartesian if and only if $(f', 1)$ is so.
Thus the following claim completely characterizes strongly cartesian
morphisms. Claim: A morphism
$$
(a, b, \alpha) :
X_1 = (U_1, \phi_1 : V_1 \to u(U_1), x_1)
\longrightarrow
(U_2, \phi_2 : V_2 \to u(U_2), x_2) = X_2
$$
of $u_{pp}\mathcal{S}$ has image $f = ((a, b, \alpha), 1)$
strongly cartesian in $u_p\mathcal{S}$ if and only if $\alpha$
is a strongly cartesian morphism of $\mathcal{S}$.

\medskip\noindent
Assume $\alpha$ strongly cartesian.
Let $X = (U, \phi : V \to u(U), x)$ be another object, and let
$f_2 : X \to X_2$ be a morphism of $u_p\mathcal{S}$
such that $p_p(f_2) = b \circ b_1$ for some $b_1 : U \to U_1$.
To show that $f$ is strongly cartesian we have to show
that there exists a unique morphism $f_1 : X \to X_1$ in $u_p\mathcal{S}$
such that $p_p(f_1) = b_1$ and $f_2 = f \circ f_1$ in $u_p\mathcal{S}$.
Write $f_2 = (f'_2 : X' \to X_2, r : X' \to X)$. Again we can write
$f_2 = (f'_2, 1) \circ (r, 1)^{-1}$ in $u_p\mathcal{S}$. Since $(r, 1)$
is an isomorphism whose image in $\mathcal{D}$ is an identity we
see that finding a morphism $f_1 : X \to X_1$ with the required
properties is the same thing as finding a morphism $f'_1 : X' \to X_1$
in $u_p\mathcal{S}$ with $p(f'_1) = b_1$ and $f'_2 = f \circ f'_1$.
Hence we may assume that $f_2$ is of the form
$f_2 = ((a_2, b_2, \alpha_2), 1)$ with $b_2 = b \circ b_1$. Here is
a picture
$$
\xymatrix{
& & (U_1, V_1 \to u(U_1), x_1) \ar[d]^{(a, b, \alpha)} \\
(U, V \to u(U), x) \ar[rr]^{(a_2, b_2, \alpha_2)} & &
(U_2, V_2 \to u(U_2), x_2)
}
$$
Now it is clear how to construct the morphism $f_1$.
Namely, set $U' = U \times_{U_2} U_1$ with projections
$c : U' \to U$ and $a_1 : U' \to U_1$.
Pick a strongly cartesian morphism $\gamma : x' \to x$
lifting the morphism $c$. Since $b_2 = b \circ b_1$,
and since $u(U') = u(U) \times_{u(U_2)} u(U_1)$ we see that
$\phi' = (\phi, \phi_1 \circ b_1) : V \to u(U')$. Since
$\alpha$ is strongly cartesian, and
$a \circ a_1 = a_2 \circ c = p(\alpha_2 \circ \gamma)$
there exists a morphism
$\alpha_1 : x' \to x_1$ lifting $a_1$ such that
$\alpha \circ \alpha_1 = \alpha_2 \circ \gamma$.
Set $X' = (U', \phi' : V \to u(U'), x')$.
Thus we see that
$$
f_1 =
((a_1, b_1, \alpha_1) : X' \to X_1, (c, \text{id}_V, \gamma) : X' \to X) :
X \longrightarrow X_1
$$
works, in fact the diagram
$$
\xymatrix{
(U', \phi' : V \to u(U'), x')
\ar[d]_{(c, \text{id}_V, \gamma)}
\ar[rr]_{(a_1, b_1, \alpha_1)}
& & (U_1, V_1 \to u(U_1), x_1) \ar[d]^{(a, b, \alpha)} \\
(U, V \to u(U), x) \ar[rr]^{(a_2, b_2, \alpha_2)} & &
(U_2, V_2 \to u(U_2), x_2)
}
$$
is commutative by construction. This proves existence.

\medskip\noindent
Next we prove uniqueness, still in the special case
$f = ((a, b, \alpha), 1)$ and $f_2 = ((a_2, b_2, \alpha_2), 1)$.
We strongly advise the reader to skip this part.
Suppose that $g_1, g'_1 : X \to X_1$ are two morphisms of
$u_p\mathcal{S}$ such that $p_p(g_1) = p_p(g'_1) = b_1$ and
$f_2 = f \circ g_1 = f \circ g'_1$. Our goal is to show that
$g_1 = g'_1$. By
Categories, Lemma \ref{categories-lemma-morphisms-right-localization}
we may represent $g_1$ and $g'_1$ as the equivalence classes of
$(f_1 : X' \to X_1, r : X' \to X)$ and
$(f'_1 : X' \to X_1, r : X' \to X)$ for some $r \in R$. By
Categories, Lemma \ref{categories-lemma-equality-morphisms-right-localization}
we see that $f_2 = f \circ g_1 = f \circ g'_1$ means that there
exists a morphism $r' : X'' \to X'$ in $u_{pp}\mathcal{S}$ such that
$r' \circ r \in R$ and
$$
(a, b, \alpha) \circ f_1 \circ r' =
(a, b, \alpha) \circ f'_1 \circ r' =
(a_2, b_2, \alpha_2) \circ r'
$$
in $u_{pp}\mathcal{S}$. Note that
now $g_1$ is represented by $(f_1 \circ r', r \circ r')$ and
similarly for $g'_1$. Hence we may assume that
$$
(a, b, \alpha) \circ f_1 =
(a, b, \alpha) \circ f'_1 =
(a_2, b_2, \alpha_2).
$$
Write $r = (c, \text{id}_V, \gamma) : (U', \phi' : V \to u(U'), x')$,
$f_1 = (a_1, b_1, \alpha_1)$, and $f'_1 = (a'_1, b_1, \alpha'_1)$.
Here we have used the condition that $p_p(g_1) = p_p(g'_1)$.
The equalities above are now equivalent to
$a \circ a_1 = a \circ a'_1 = a_2 \circ c$ and
$\alpha \circ \alpha_1 = \alpha \circ \alpha'_1 = \alpha_2 \circ \gamma$.
It need not be the case that $a_1 = a'_1$ in this situation.
Thus we have to precompose by one more morphism from $R$.
Namely, let $U'' = \text{Eq}(a_1, a'_1)$ be the equalizer of
$a_1$ and $a'_1$ which is a subobject of $U'$. Denote
$c' : U'' \to U'$ the canonical monomorphism.
Because of the relations among the morphisms
above we see that $V \to u(U')$ maps into
$u(U'') = u(\text{Eq}(a_1, a'_1)) = \text{Eq}(u(a_1), u(a'_1))$.
Hence we get a new object $(U'', \phi'' : V \to u(U''), x'')$, where
$\gamma' : x'' \to x'$ is a strongly cartesian morphism lifting $\gamma$.
Then we see that we may precompose $f_1$ and $f'_1$ with the element
$(c', \text{id}_V, \gamma')$ of $R$. After doing this, i.e., replacing
$(U', \phi' : V \to u(U'), x')$ with
$(U'', \phi'' : V \to u(U''), x'')$, we get back to the previous
situation where in addition we now have that $a_1 = a'_1$.
In this case it follows formally from the fact that $\alpha$
is strongly cartesian (!) that $\alpha_1 = \alpha'_1$.
This shows that $g_1 = g'_1$ as desired.

\medskip\noindent
We omit the proof of the fact that for any strongly cartesian morphism
of $u_p\mathcal{S}$ of the form $((a, b, \alpha), 1)$ the morphism
$\alpha$ is strongly cartesian in $\mathcal{S}$.
(We do not need the characterization of strongly cartesian morphisms
in the rest of the proof, although we do use it later in this section.)

\medskip\noindent
Let $(U, \phi : V \to u(U), x)$ be an object of $u_p\mathcal{S}$.
Let $b : V' \to V$ be a morphism of $\mathcal{D}$. Then the morphism
$$
(\text{id}_U, b, \text{id}_x) :
(U, \phi \circ b : V' \to u(U), x)
\longrightarrow
(U, \phi : V \to u(U), x)
$$
is strongly cartesian by the result of the preceding paragraphs and we win.
\end{proof}

\end{document}
